% -*- mode: latex -*- %

\begin{abstract}
  The construction of most supervised learning datasets revolves around
  collecting multiple labels for each instance, then aggregating the labels
  to form a type of ``true'' label.
  %
  We question the wisdom of this pipeline by developing a (stylized)
  theoretical model of this process and analyzing its statistical
  consequences, showing how access to non-aggregated label information can
  make training well-calibrated models more feasible than it is with
  cleaned labels.
  %
  The entire story, however, is subtle, and the
  contrasts between aggregated and fuller label information depend on
  the particulars of the problem,
  where estimators that use aggregated information exhibit robust but slower
  rates of convergence, while estimators that can effectively
  leverage all labels converge more quickly \emph{if} they have
  fidelity to (or can learn) the true labeling process.
  The theory
  makes several predictions for real-world datasets, including when
  non-aggregate labels should improve learning performance, which we
  test to corroborate the validity of our predictions.
\end{abstract}
